\section{Discussion}

\textbf{Why rules help.} Most of the gain is in refusing correctly: on the first test set, 87 more correct refusals against one fewer
correct \emph{answer} decision. Out-of-scope questions in real traffic fall into a few recurring kinds (file status, complaints and
petitions, penalties, police internal matters) whose wording is regular enough for patterns and whose cost of error is high. The learned
components answer a different question. The reranker scores how well a procedure matches a text, which is high for a complaint that
mentions a real procedure; the LLM reader recognizes intents but often labels long narrative procedural questions as complaints or legal
questions. Its procedure name alone gave the best first-test scores, but it lowered end-to-end accuracy on development in both rounds, so
we do not claim that it helps.

\textbf{Remaining errors of A.} In a post-hoc analysis of the 24 \emph{answer} questions that A misses on the first test set, 12 are
refused; for 4 of them the base system had ranked the correct procedure first, so the out-of-scope gate is at fault. Of the 12 that A
answers wrongly or clarifies, 7 have the correct procedure outside its top five: long letters in which the procedure is implied rather
than named (a car registered to a deceased parent, while the story triggers death registration). A2 was written against such errors.

\textbf{Why the second round did not carry over.} A2 + B changes the decision or the chosen procedure of about as many questions on the new
set as on development (5.2\% against 6.2\%), but, counting a decision as correct only with the right procedure, the changes were 30 gains
and no losses on development and 4 and 4 on the new set. The clearest sign is in refusals: its new out-of-scope patterns turned 15 wrong
answers into correct refusals on development but one of the 335 questions to be refused in the new set, which suggests that they were
written from particular phrasings. Two other factors may contribute: with 21 answerable questions the new set has little power on the
\emph{answer} side, and its labels were proposed by a different model with a slightly extended guideline. Either way, a gain of 25
questions on reused questions became a tie on new ones.

\textbf{Limitations and threats to validity.} (i) One domain, language and data snapshot; the rules are hand-written for Vietnamese
commune-level procedures, and the results say nothing about provincial or ministry procedures or a general assistant for all procedures. (ii) The test
sets have only 59 and 21 \emph{answer} and 6 and 7 \emph{clarify} questions, so small differences are not significant and nothing can be
concluded about clarification. (iii) The development and first test labels were produced with the same AI assistant that wrote the rules;
the second-pass sample and the second test set were labeled with another model and checked item by item by an author, who had seen the
aggregate class proportions of the first test set. Agreement is substantial and conclusions hold under either label set, but no fully
independent human annotation exists. (iv) All questions come from the same two portals, and 79\% refusals in our samples understate the
share in random traffic. (v) Systems from other studies could not be run on our questions, and their published numbers measure different
tasks; our external baseline is a RAG design we implemented. (vi) LLM components were run once at temperature 0 on one quantized model.
